\chapter[Oscillators and stochastic dynamics]{Oscillators, bifurcations, and stochastic dynamics}
\label{ch:oscillator-foundations}
Oscillatory-looking trajectories can emerge from a dynamical
system, an externally driven filter, or chance structure in
a finite recording. These are mathematically distinct accounts.

\section{Hopf normal form and an actual periodic attractor}
The complex-valued ordinary differential equation
$\dot z=(\mu+i\omega)z-|z|^2z$ with $\omega\ne0$
can be written $z=re^{i\theta}$, giving
\[
\dot r=\mu r-r^3,\qquad \dot\theta=\omega.
\]
For $\mu<0$ the origin attracts nearby trajectories.
For $\mu>0$ the origin repels, while the circle
$r=\sqrt\mu$ is an attracting periodic orbit. Its period
is $2\pi/|\omega|$.
This explicit derivation demonstrates a supercritical Hopf-type
normal form; deciding whether experimental LFP oscillations
arise from an actual Hopf bifurcation requires more evidence.

\section{Phase reduction and weak coupling}
A stable limit-cycle oscillator with phase coordinate $\theta$
responds to weak perturbation $u(t)$ to leading order as
\[
\dot\theta=\omega+Z(\theta)\cdot u(t),
\]
where $Z$ is the infinitesimal phase-response curve (PRC).
For two weakly coupled oscillators, averaging can produce
a phase-difference equation
$\dot\psi=\Delta\omega+\varepsilon H(\psi)$.
Locking requires a stable fixed point of this equation;
oscillatory phase locking need not imply anatomical
synapses, and non-sinusoidal harmonics can generate apparent PAC.

\section{jPCA and skew-symmetric dynamics}
If $\dot x=Mx$ for a real skew-symmetric matrix
$M^\top=-M$, then
\[
\frac{d}{dt}\|x\|_2^2=2x^\top Mx=0.
\]
A real skew-symmetric matrix has orthogonal real normal form
consisting of planar blocks
$\begin{psmallmatrix}0&-\omega_j\\\omega_j&0\end{psmallmatrix}$
(and possibly zeros); its nonzero eigenvalues occur as
imaginary pairs $\pm i\omega_j$. Thus a fitted rotational
component conserves Euclidean energy in the ideal autonomous
model. The file \path{jpca_dynamics.m}
fits such a skew generator to finite-difference derivatives,
not a complete state-space model or a proof of physiological cycles.

\section{Ornstein--Uhlenbeck noise and the Fokker--Planck equation}
For an It\^o stochastic differential equation
$dX_t=a(X_t)dt+b(X_t)dW_t$, a twice differentiable
function $f$ obeys It\^o's formula
\[
df(X_t)=\left(af'+\tfrac12 b^2f''\right)dt
+b f'\,dW_t.
\]
Its density, under regularity and boundary assumptions,
evolves by the Fokker--Planck equation
$\partial_tp=-\partial_x(ap)
 +\tfrac12\partial_{xx}(b^2p)$.
For the Ornstein--Uhlenbeck process
$dX_t=-\kappa X_tdt+\sigma dW_t$, with $\kappa>0$,
the stationary variance is $\sigma^2/(2\kappa)$ and the
stationary autocovariance is
$\sigma^2e^{-\kappa|\tau|}/(2\kappa)$.
An exponential autocorrelation is therefore compatible with
a stochastic relaxation process and need not imply an intrinsic
neuronal pacemaker or periodic oscillator.

\begin{example}[Simulation check]
Simulate a damped OU fluctuation and a noisy Hopf limit cycle
at the same sampling frequency. Compare autocorrelation,
PSD and phase stability under different recording lengths.
A narrow spectral peak alone does not establish the Hopf model.
\end{example}